\documentclass[11pt,a4paper]{article}
\usepackage[utf8]{inputenc}
\usepackage[T1]{fontenc}
\usepackage[french]{babel}
\usepackage{geometry}
\usepackage{hyperref}
\usepackage{enumitem}

\geometry{margin=2.5cm}

\title{\textbf{Innovative Student Project}\\Project Definition and Brainstorming}
\author{Group 3  IDRIS, AHMED , SALIMA }
\date{}

\begin{document}

\maketitle

\section*{Project Overview}
Our project is an innovative web-based application designed to help students better manage their time, mental well-being, and academic organization. It combines planning, self-tracking, and personalized feedback in one accessible tool.

\section*{1. What do we want to achieve through our innovative project?}
We want to create a simple but intelligent student assistant that improves existing tools by making them more personalized, more accessible, and specifically adapted to student life.  
Rather than inventing a new technology, our innovation lies in improving usability and relevance.

\section*{2. What do we want others to understand through our work?}
We chose this project because we personally experienced academic stress, irregular routines, and difficulty managing workload.  
Through this project, we want others to understand that student well-being is a real issue and that technology can support students when designed with their needs in mind.

\section*{3. What specific problem are we trying to solve?}
Students today face:
\begin{itemize}
  \item High academic pressure
  \item Stress and anxiety
  \item Poor time management
  \item Irregular sleep schedules
  \item Risk of burnout
\end{itemize}
Existing tools are often fragmented, too complex, or not adapted to students.

\section*{4. What proof do we have that this problem exists?}
Several studies and reports (WHO, universities, academic research) show:
\begin{itemize}
  \item A significant increase in stress and anxiety among students
  \item More than 50\% of university students report high levels of stress
  \item Sleep irregularity negatively affects learning and memory
\end{itemize}
This confirms that the problem is real and documented.

\section*{5. What will our outcome or product look like?}
Our product will be a web-based application including:
\begin{itemize}
  \item A weekly planning interface
  \item Mood and stress tracking
  \item Sleep and habit monitoring
  \item A dashboard with visual graphs
  \item Personalized suggestions
\end{itemize}

\section*{6. What type of materials will we use?}
\begin{itemize}
  \item Computers for development
  \item Web browsers for usage
  \item Free development tools (HTML, CSS, JavaScript)
  \item GitHub for collaboration
  \item Free design tools (Figma / Canva)
\end{itemize}

\section*{7. What techniques / approaches will we use?}
\begin{itemize}
  \item User-centered design
  \item Iterative development (step-by-step improvement)
  \item Simple data visualization
  \item Basic personalization logic
\end{itemize}

\section*{8. What type of information should we include?}
\begin{itemize}
  \item User-input data (mood, sleep, schedule)
  \item Educational tips about productivity and well-being
  \item Weekly summaries
  \item Visual indicators of progress
\end{itemize}

\section*{9. How will we present the information?}
Information will be presented through:
\begin{itemize}
  \item A clear and simple dashboard
  \item Graphs and visual summaries
  \item Short and understandable messages
  \item Clean interface design
\end{itemize}

\section*{10. Should we include visuals?}
Yes. We will include:
\begin{itemize}
  \item A visual prototype
  \item A poster summarizing the project
  \item Possibly a short demonstration video
\end{itemize}

\section*{11. Do we need to consider intellectual property issues?}
Yes. We will ensure:
\begin{itemize}
  \item The project idea and design are original
  \item No direct copying of existing applications
  \item Only free or open-source resources are used
\end{itemize}

\section*{12. Who is the audience (target market)?}
\begin{itemize}
  \item Primary audience: university students and high school students
  \item Secondary audience: teachers and institutions interested in student support tools
\end{itemize}

\section*{13. How will we get feedback about the success of the project?}
We will:
\begin{itemize}
  \item Ask students to test the prototype
  \item Use questionnaires (Google Forms)
  \item Collect feedback about usability, usefulness, and clarity
  \item Improve the project based on this feedback
\end{itemize}

\end{document}
